\section{Related Work}
\label{sec:related}

\paragraph{Sensitivity and diversity of multi-task rankings.}
\citet{zhang2024diversity} treat a multi-task benchmark as a social-choice problem, separate cardinal from ordinal aggregation, and show a trade-off between the diversity of a benchmark and the stability of its ranking when noise is injected into tasks. That perturbation is untargeted. We ask what a developer who \emph{chooses} the tasks must spend to put one given model first, and use their cardinal/ordinal taxonomy and the dispersion and correlation of task scores as candidate predictors of the price (\Cref{sec:predictors}). Finite-sample robustness metrics ask how little data must change to flip a conclusion \citep{broderick2020dropping}; ours is the smallest change to one model's scores that flips a leaderboard.

\paragraph{Strategic post-training and contamination.}
\citet{chen2026incentives} model a designer and several developers who each choose a scalar benchmark-specific effort; their benchmark has one score, so the aggregation across tasks, our subject, is absent, and their results are theoretical. We give the single-developer, multi-task measurement on real leaderboards. \citet{dominguez2024testtask} discuss training on data related to the evaluation tasks as a confound, and \citet{zhang2025trainbeforetest} propose giving every model the same task-specific fine-tuning before testing; we cite both for the mechanism. Private testing on Chatbot Arena \citep{singh2025illusion}, benchmark data in training \citep{zhou2023cheater} and rephrased test samples \citep{yang2023rephrased} show that the actions we price occur; they do not measure a price per aggregation rule.

\paragraph{Aggregation rules and social choice.}
Voting rules have been applied to rank agents and models over tasks \citep{lanctot2023social,rofin2023votenrank}, Nash averaging makes the ranking invariant to redundant tasks and models \citep{balduzzi2018reevaluating}, and the choice of tasks can itself reorder methods \citep{dehghani2021lottery}; these works compare rules by axioms or ranking agreement, not by the cost of moving one model to the top. Our price is a cardinal bribery problem with ceilings \citep{brandt2016comsoc,bredereck2016shift}, and the intuition from robust statistics is that a rule with a high breakdown point, such as the median, is expensive to capture. For a finite margin of victory the price is instead a function of the target's deficits, and the intuition has the wrong sign on HELM and depends on board and budget elsewhere (\Cref{sec:pricemap,sec:breadth}).
